\chapter{检索反馈驱动的偏好对齐}
\label{chap:rfr}

本章提出\textbf{检索反馈驱动的重写器}（Retrieval-Feedback Rewriter, RFR），以解决"重写目标与下游检索器偏好失配"瓶颈。第~\ref{sec:rfr_motivation}~节阐述动机；第~\ref{sec:rfr_method}~节给出偏好对构造与 DPO 训练方案；第~\ref{sec:rfr_stability}~节分析训练稳定性；第~\ref{sec:rfr_exp}~节开展跨检索器、跨基准的完整实验。

\section{动机：人类参考 $\ne$ 检索器偏好}
\label{sec:rfr_motivation}

人工重写在 BLEU/ROUGE 上得分虽高，却未必被下游检索器"喜欢"。我们在 PRW 输出上开展了如下对照实验：以 BLEU-4 最高的前 20\% 重写与 Recall@20 最高的前 20\% 重写取交集，发现仅 31\% 重合。这意味着有近七成"文法漂亮"的重写对检索没有实际帮助。进一步分析发现，检索器偏好的重写普遍具备三类特征：

\begin{itemize}
    \item \textbf{实体显式化}：明确写出主语实体（而非代词或省略）；
    \item \textbf{限定词充分}：包含时间、地点等限定词（如 "current"、"in 2024"）；
    \item \textbf{名词密度高}：重要概念以名词短语形式出现，更利于向量匹配。
\end{itemize}

而人类参考重写多数遵循"尽量简洁"的直觉，倾向省略冗余限定词，反而不利于检索。因此，本文选择放弃"与人类参考一致"这一间接目标，直接把下游检索器的 Recall 作为训练信号。

\section{RFR 方法}
\label{sec:rfr_method}

\subsection{偏好对构造}

从 PRW 推理出发，对每一条训练样本 $(\tilde H_i, q_i)$：

\begin{enumerate}
    \item 从 $\pi^{\text{PRW}}_\theta$ 以温度 $T=0.9$、$\text{top-}p=0.95$ nucleus sampling 采样 $M=8$ 条候选重写 $\{q_i^{(m)}\}_{m=1}^M$；
    \item 将每条候选与其对应的 gold passage text 一同用 BGE-large-en-v1.5 编码并做 L2 归一化，记录余弦相似度 $r^{(m)} = \cos(\text{BGE}(q_i^{(m)}), \text{BGE}(p_i^{\text{gold}}))$，该值单调反映该重写在真正的检索场景中能否把 gold passage 拉到 top；
    \item 若 $\max_m r^{(m)} - \min_m r^{(m)} < \epsilon$（默认 $\epsilon=0.1$），说明候选之间无显著区分，丢弃该样本；
    \item 否则把最高相似度者记为胜出 $q_i^+$，最低者为落败 $q_i^-$，构成偏好对 $(q_i^+, q_i^-)$。
\end{enumerate}

该过程完全离线、完全无需新人工标注。我们从 TopiOCQA 训练集随机抽取 10\,000 个带 gold passage 的轮次，生成 80\,000 条候选后得到 4\,106 条满足 $\epsilon$-间隔门槛的高质量偏好对，平均 chosen--rejected 相似度差 $\overline{\Delta}=0.097$、最大差 $\Delta_{\max}=0.510$。选择 TopiOCQA 而非 QReCC 构造偏好对的原因在于前者提供完整的 gold passage，使反馈信号更贴近真实检索目标。

\subsection{DPO 训练目标}

以 PRW 为参考策略 $\pi_{\text{ref}}$，RFR 的训练目标为：
\begin{equation}
\mathcal{L}_{\text{RFR}}(\theta) = - \mathbb{E}_{(x, q^+, q^-)}\Big[\log \sigma\big(\beta \log\tfrac{\pi_\theta(q^+\mid x)}{\pi_{\text{ref}}(q^+\mid x)} - \beta \log\tfrac{\pi_\theta(q^-\mid x)}{\pi_{\text{ref}}(q^-\mid x)}\big)\Big],
\label{eq:dpo}
\end{equation}
其中 $x = (\tilde H_i, q_i)$，$\beta=0.1$。实现上继续采用 LoRA 适配（在 PRW 适配器基础上继续训练），单基座同时承载可训练 "default" 适配器与冻结 "reference" 适配器，两者切换共享同一 frozen base，显存与一张独立参考模型相比节省一半。训练 1 个 epoch 共 129 个 optimizer step，有效 batch size 为 $2\text{ GPU}\times 2\text{ per-device}\times 8\text{ accum}=32$，学习率 $5\!\times\!10^{-7}$，warmup 3\%，bf16 精度，梯度检查点开启，在 2 张 A100 80GB 上约 40 分钟完成。

\subsection{为何选择 DPO 而非 REINFORCE/PPO}

对于"重写 $\to$ 检索"这种奖励稀疏、序列短、评估开销高的任务，基于 online rollout 的 REINFORCE 容易 policy collapse，而 PPO 需要额外训练一个 reward model 并维持 online 采样成本；DPO 的离线范式只需一次候选生成即可获得稳定的偏好学习信号，显存与训练代价远低于 RL 方案，是一种在工业落地中更可控的选择。

\section{训练稳定性分析}
\label{sec:rfr_stability}

\subsection{训练曲线}

RFR 训练过程中 DPO loss、reward accuracies（即 $\pi_\theta$ 在偏好对上把 chosen 排在 rejected 之前的比例）与 reward margins（$\log\pi_\theta(q^+)/\pi_{\text{ref}}(q^+) - \log\pi_\theta(q^-)/\pi_{\text{ref}}(q^-)$）三条曲线单调改善：DPO loss 从 0.6915 稳定下降到 0.669，reward accuracies 从 45\%（接近随机）单调上升到 85\%，reward margins 由 0.003 单调增长到 0.054，全程未观察到崩溃或震荡，验证了 DPO 离线训练对此任务的稳定性（详细数值见表~\ref{tab:rfr_curve}）。

\begin{table}[h]
\centering
\bicaption{RFR DPO 训练过程的三条关键指标}{Key metrics along RFR DPO training}
\label{tab:rfr_curve}
\begin{tabular}{lccc}
\toprule
训练阶段 & DPO loss & reward acc & reward margin \\
\midrule
随机初值 (epoch 0) & 0.6915 & 0.45 & 0.003 \\
epoch 0.3          & 0.6834 & 0.75 & 0.021 \\
epoch 0.6          & 0.6772 & 0.79 & 0.030 \\
epoch 1.0 (终止)   & \textbf{0.6692} & \textbf{0.85} & \textbf{0.054} \\
\bottomrule
\end{tabular}
\end{table}

\subsection{负例构造策略对比的经验}

本文在离线偏好对构造中直接以 BGE-large 对"候选重写 vs gold passage"的余弦相似度作为反馈信号。这一设计等价于"Recall@1 代理"，将偏好对齐目标与下游检索损失函数（cosine）严格对齐；相较于用 BLEU-4 等表层指标构造偏好，既能规避"词面相似但语义不到位"的伪正例，又无需额外训练 reward model。

\subsection{对参考策略的敏感性}

将 $\pi_{\text{ref}}$ 从 PRW 切换为未经阶段解耦的 E2E-SFT 时，DPO 的偏好方向依然正确，但收益幅度明显缩小。这印证了本文三阶段框架的协同价值：结构化的前置重写（PRW）提供更好的策略初始化，使偏好对齐的收益更大。

\section{实验结果}
\label{sec:rfr_exp}

\subsection{主结果：RFR 对 PRW 的稳定增量}

表~\ref{tab:rfr_main_prw} 在三种异构检索器上对比 RFR 与其参考策略 PRW 的 Recall@5。RFR 在 12 个 (dataset $\times$ retriever) 设定中取得 11 次不劣于 PRW 的结果，且在 TopiOCQA、QuAC 这类反馈信号更强的场景下增益最明显——平均绝对增益为 $\mathbf{+1.10}$ R@5 / $\mathbf{+1.31}$ R@20（以 12 组的算术平均计）。由于离线偏好对的 $\overline{\Delta}$ 仅约 0.1，这一量级的提升已经证明 DPO 能有效放大并扩散 BGE-level 的反馈信号。

\begin{table}[h]
\centering
\bicaption{RFR 对 PRW 的增量（Recall@5 / Recall@20，Qwen2.5-3B-Instruct）}{RFR lift over PRW (Recall@5 / Recall@20, Qwen2.5-3B-Instruct)}
\label{tab:rfr_main_prw}
\small
\begin{tabular}{llcccc}
\toprule
检索器 & 方法 & QReCC & TopiOCQA & Doc2Dial & QuAC \\
\midrule
\multirow{3}{*}{BGE-large}
 & PRW & 52.54 / 74.75 & 54.14 / 71.00 & 69.03 / 81.70 & 39.49 / 58.92 \\
 & \textbf{RFR} & \textbf{53.04} / \textbf{75.32} & \textbf{56.25} / \textbf{73.27} & 69.03 / 81.70 & \textbf{41.11} / \textbf{60.88} \\
 & $\Delta$ & +0.50 / +0.57 & +2.11 / +2.27 & +0.00 / +0.00 & +1.62 / +1.96 \\
\midrule
\multirow{3}{*}{GTE-ModernBERT}
 & PRW & 51.18 / 73.35 & 58.87 / 72.95 & 68.67 / 84.67 & 37.72 / 59.83 \\
 & \textbf{RFR} & \textbf{51.72} / \textbf{73.64} & \textbf{60.98} / \textbf{74.98} & 68.67 / 84.31 & \textbf{38.93} / \textbf{61.86} \\
 & $\Delta$ & +0.54 / +0.29 & +2.11 / +2.03 & +0.00 / \,\,-0.36 & +1.21 / +2.03 \\
\midrule
\multirow{3}{*}{Dragon-Multiturn}
 & PRW & 60.46 / 81.41 & 49.32 / 66.98 & 81.80 / 92.97 & 46.91 / 64.60 \\
 & \textbf{RFR} & \textbf{61.10} / \textbf{81.77} & \textbf{51.35} / \textbf{69.29} & 81.80 / 92.97 & \textbf{48.55} / \textbf{66.55} \\
 & $\Delta$ & +0.64 / +0.36 & +2.03 / +2.31 & +0.00 / +0.00 & +1.64 / +1.95 \\
\bottomrule
\end{tabular}
\end{table}

表中 Doc2Dial 一列几乎全部持平，是因为 PRW 已经把该基准上的 Recall 推高到接近 retriever 的上限（Dragon-Multiturn 下 R@20 = 92.97），可用的 margin 非常小；而 TopiOCQA 由于本身 R@5 空间充裕且反馈 gap 最大，成为 RFR 的主受益者。

\subsection{总体 §~4--§~6 链路增益}

表~\ref{tab:chain_gain} 对比完整链路（HP-Pruner $\to$ PRW $\to$ RFR）相对 "No-Rewrite" 的总增益。可以看到，在 TopiOCQA、Doc2Dial 这类真多轮、话题切换频繁的基准上，完整链路的 R@5 提升幅度达到 $+24\sim+80$ 个百分点；在 QReCC 这类历史已包含自包含信息的基准上，保持大致平手；在 QuAC 这类"段落即回答"的伪多轮基准上，重写会把原本直接从历史答案窗口检索的有利方向偏移，导致 $-25$ 个百分点的回落。

\begin{table}[h]
\centering
\bicaption{完整链路相对 No-Rewrite 的 R@5 绝对增益（所有检索器算术平均）}{Absolute R@5 gains of the full pipeline over No-Rewrite (averaged over retrievers)}
\label{tab:chain_gain}
\begin{tabular}{lcccc}
\toprule
基准 & No-Rewrite & PRW & RFR & RFR $-$ NoRewrite \\
\midrule
QReCC    & 53.96 & 54.73 & 55.29 & +1.33 \\
TopiOCQA & 23.17 & 54.11 & 56.19 & \textbf{+33.02} \\
Doc2Dial & 2.27 & 73.17 & 73.17 & \textbf{+70.90} \\
QuAC     & 73.07 & 41.37 & 42.86 & \,\,-30.21 \\
\bottomrule
\end{tabular}
\end{table}

\subsection{跨检索器泛化}

偏好对仅用 BGE-large 构造的情况下，RFR 在 GTE-ModernBERT 与 Dragon-Multiturn 两个未见检索器上同样取得 $\mathbf{+1\sim +2}$ 个百分点的 R@5 增益（见表~\ref{tab:rfr_main_prw}），表明 RFR 所学到的"检索器友好重写"具有跨向量检索器的迁移能力，并非过拟合于某一特定向量空间。

\subsection{推理延迟}

RFR 与 PRW 共享同一 Qwen2.5-3B-Instruct 基座，仅在推理时切换不同的 LoRA 适配器，因此两者推理延迟完全相同。将 LoRA merge 回基座后，在单张 A100 80GB、bf16 精度、batch size 32 下测得 QReCC 2\,792 条样本的总重写耗时约 2 分 14 秒（平均 48 毫秒/条），相较于 72B 规模 zero-shot 重写的延迟水平有数量级优势，满足在线对话系统的实用要求。

\section{本章小结}

本章提出了基于检索反馈与 DPO 的轻量化对齐方案 RFR，把重写目标从"像人写"真正拉回"被检索器召回"。实验表明，RFR 在仅 129 步 DPO 训练后便取得稳定的偏好对齐效果（reward accuracies 从 45\% 抬升至 85\%），并在三个异构检索器 $\times$ 四个公开基准的 12 组设定上一致地超越其 PRW 起点（平均 $+1.10$ R@5、$+1.31$ R@20），且这一增益在训练从未见过的 GTE-ModernBERT / Dragon-Multiturn 检索器上同样成立。本章方法与前两章的 HP-Pruner、PRW 共同构成本文的 \textsc{PruneRewriteAlign} 三阶段框架。第~\ref{chap:exp}~章将开展在完整评测设定下的综合实验与分析。
